% SOURCE_STATEMENT: "Develop a mathematical framework that axiomatizes physics, particularly mechanics, thermodynamics, and probability theory."
% Preparation and final self-audit: CONTRIBUTING.md prompts C and D; source/audit artifacts retained in the local physics_statements_2026_10_11 batch.
% !TeX program = lualatex
\documentclass[11pt]{article}
\usepackage[margin=25mm]{geometry}
\usepackage{fontspec}
\setmainfont{Latin Modern Roman}
\usepackage{amsmath,amssymb,mathtools}
\usepackage{unicode-math}
\setmathfont{Latin Modern Math}
\usepackage{xurl,hyperref}
\hypersetup{colorlinks=true,urlcolor=blue}
\setcounter{secnumdepth}{0}
\setlength{\parindent}{0pt}
\setlength{\parskip}{5pt}
\setlength{\emergencystretch}{3em}
\newcommand{\sref}[2]{\hyperlink{source:#1:#2}{[S#2]}}
\newcommand{\cataloguescope}[1]{\paragraph{Recorded target.}#1}
\newcommand{\R}{\mathbb R}
\newcommand{\C}{\mathbb C}
\newcommand{\N}{\mathbb N}
\newcommand{\Z}{\mathbb Z}
\newcommand{\Q}{\mathbb Q}
\newcommand{\F}{\mathbb F}
\newcommand{\eps}{\varepsilon}
\begin{document}
% UnsolvedMath ID 61; source catalogue code HIL-006
\setcounter{section}{60}
\section{Hilbert's 6th Problem: Axiomatization of Physics}\label{61}
{\small UnsolvedMath ID 61 · Mathematical Physics}
\subsection{Definitions and mathematical statement}

\paragraph{Shared notation.}$\mathbb N=\{1,2,\ldots\}$; the standard number fields have their usual meanings. All additional conventions are stated locally.


\paragraph{Definitions and notation.}An axiomatization specifies mathematical objects, axioms governing their relations and evolution, and the deductions to be established within that framework. In mechanics, a microscopic model describes individual particles; a continuum model describes fields such as mass density and velocity. A limit from one to the other must specify a scaling, topology of convergence and admissible initial data. Hilbert’s original sixth problem is a programme concerning axioms of physical theories, especially probability and mechanics; it is not a single formal proposition.

\paragraph{Catalogue wording retained for traceability.}
Develop a mathematical framework that axiomatizes physics, particularly mechanics, thermodynamics, and probability theory. Hilbert’s original programme especially calls for an axiomatic treatment of probability, including mean-value methods in mathematical physics, and a rigorous development of limiting processes leading from atomic mechanics to continuum mechanics.
\paragraph{Scope and source qualification.}The catalogue includes thermodynamics as a later framing. Hilbert supplies no fixed final axiom system or universally agreed completion criterion. Existing axiomatizations of particular theories constitute context rather than a settlement of this entire programme.
\subsection{Short English statement}
Give rigorous axioms for physical theories and justify the passage from microscopic particle mechanics to continuum descriptions.
\subsection{Sources}
\begin{itemize}
\item[\hypertarget{source:61:1}{S1}] ULAM.AI and the UnsolvedMath Contributors, UnsolvedMath, record 61 (HIL-006). Catalogue curation under CC BY 4.0; definitions and English exposition added locally.
\url{https://huggingface.co/datasets/ulamai/UnsolvedMath}.
\item[\hypertarget{source:61:2}{S2}] David Hilbert, Mathematical Problems (1900), Problem 6, English translation archived by D. E. Joyce.
\url{https://mathcs.clarku.edu/~djoyce/hilbert/problems.html}.
\end{itemize}
\label{end:61}
\end{document}
